\chapter{Dossiers commentés}\label{ch:dossiers}

Ce chapitre réunit des dossiers sur quarante-sept textes du corpus, choisis pour leur intérêt philologique, chronologique ou historique, et en particulier tous les textes dont dépendent les conclusions des chapitres précédents. Chaque dossier donne, après une ligne de données (année de classement et fondement, références, rattachement topographique, support selon l’EDB, renvoi à la notice du catalogue), le texte transmis par l’EDCS, ligne à ligne, puis le texte de l’EDB lorsqu’il existe. Les deux témoins sont reproduits sans harmonisation : leurs différences de notation (⸢ ⸣ et \emph{(!)} pour l’EDCS ; ⟨ ⟩, ⌜ ⌝, lettres soulignées ou pointées et descriptions de signes entre doubles parenthèses pour l’EDB) sont celles des bases (chapitre~\ref{ch:edition}). Viennent ensuite une traduction et un commentaire.

Les traductions sont littérales ; les crochets y signalent les passages qui traduisent une restitution. Elles ne remplacent pas une traduction établie sur l’édition imprimée. Aucun des textes n’a été collationné sur la pierre ou sur une photographie, à l’exception partielle du \nnum{38}.

\dossier{4}{Iunianus et Victora, « amis des pauvres »}
\trad{Déposition de Iunianus, la veille des ides d’avril (12 avril), sous le consulat de Marcellinus et de Probinus. Il vécut quarante ans et mourut en paix ; ami des pauvres, il vécut avec sa \emph{virginia} quinze ans. À lui, bien méritant, sa \emph{virginia} Victora (\emph{Bictora}), bien méritante, a fait (ce monument) : amie des pauvres et travailleuse.}
\comm{La plus ancienne épitaphe du corpus qui emploie \emph{virginia} pour l’épouse épousée vierge (chapitre~\ref{ch:personnes}). Les deux époux reçoivent la même qualification d’« amis des pauvres » (\emph{amator}, \emph{amatrix pauperum}, avec la graphie \emph{paupe\{ro\}rum} que l’EDCS corrige et que l’EDB conserve) ; l’épouse est en outre \emph{operaria}, mot qui peut désigner une travailleuse manuelle ou, dans un sens plus particulier, une personne active dans les œuvres charitables, ce que le texte ne permet pas de préciser. Le bêtacisme est constant (\emph{bixit}, \emph{birginia}, \emph{Bictora}). L’EDB décrit un monogramme du Christ à la fin du texte, que l’EDCS ne signale pas. La date de 341 repose sur une formule consulaire lue en entier.}

\dossier{15}{Cassia, \emph{clarissima femina}}
\trad{[…] Cassia, (fille ou épouse) de Piso, \emph{clarissima femina}, vécut trente ans et quatre-vingt-treize jours ; née sous le consulat de Sabinus et Rufinus, le dix-septième jour avant les calendes de juillet (15 juin 316), elle rendit (l’âme) après le consulat d’Amantius et d’Albinus, le dix-septième (EDB : seizième) jour avant les calendes d’octobre (15 ou 16 septembre 346).}
\comm{Le génitif \emph{Pisonis} peut désigner le père ou le mari. Le texte est gravé sur un sarcophage (17 × 130 cm selon l’EDB, qui donne sans doute les dimensions du seul champ inscrit). Il contient deux dates consulaires, celle de la naissance (316) et celle de la mort (post-consulat de 345, soit 346) ; l’année de classement est celle de la mort. La durée de vie gravée, trente ans et quatre-vingt-treize jours, concorde avec les deux dates, selon le mode de décompte (chapitre~\ref{ch:calendrier}) : du 15 juin au 15 septembre, on compte quatre-vingt-treize jours en décompte inclusif, et autant jusqu’au 16 septembre en décompte exclusif. Les deux bases divergent sur le quantième du décès (XVII ou XVI), et l’EDB restitue le nom du mois ; la concordance arithmétique ne permet pas de choisir entre les deux lectures. Le verbe \emph{reddit}, employé absolument, appartient au vocabulaire de la mort comme restitution (de l’âme ou de la vie).}

\dossier{22}{Caianus, juriste}
\trad{(Vers) Heureuse fut la vie de cet homme, plus heureux encore son trépas : le nom de Caianus croît sans cesse au fil des siècles ; la renommée ne sait pas mourir, mais vit toujours et partout. Hôte romain, il vint dans la ville du prince ; il fut ici d’abord jurisconsulte et ami. — Il repose en paix ; déposé le quatrième jour avant les nones d’août (2 août), sous le consulat de Flavius Philippus et de Flavius Salia (348). Son père Sabbatius a fait (ce monument).}
\comm{Épitaphe métrique (CLE 659) d’un jurisconsulte. Les quatre hexamètres reprennent les lieux communs de la gloire posthume (\emph{nescit fama mori}) ; la partie en prose emploie le formulaire chrétien (\emph{quiescit in pace}, avec la graphie \emph{pacem}). Le vers \emph{advenit hospes Romanus princeps in urbem} est obscur : \emph{princeps} peut qualifier la ville (« la ville capitale ») ou désigner le prince ; la traduction proposée reste incertaine. L’EDB décrit un oiseau et un chrisme radié. Le père porte un nom dérivé du sabbat, \emph{Sabbatius} (chapitre~\ref{ch:noms}).}

\dossier{26}{Esuvia Ianuaria et Aradius Melissus}
\trad{D’une bonté et d’une sainteté admirables : à Esuvia Ianuaria, son épouse, qui vécut avec moi dix-huit ans, cinq mois et vingt-cinq jours, Aradius Melissus, son mari, a fait (ce monument), et pour lui-même ; déposée en paix le troisième jour avant les calendes d’avril (30 mars), sous le consulat de Philippus et de Salia (348).}
\comm{Le texte ne donne pas l’âge de la défunte, mais la durée de la vie commune (\emph{vixit mecum}), avec la même précision qu’un âge ; la table des personnes ne lui attribue donc pas d’âge (chapitre~\ref{ch:personnes}). Le mari prépare sa propre sépulture (\emph{et sibi fecit}). Le texte n’est pas apparié à l’EDB ; il est connu par l’ICVR II 5711 et par une publication de 2012.}

\dossier{36}{Clodius Adelfius, ancien préfet de la Ville}
\trad{Clodius Adelfius, \emph{vir clarissimus}, ancien préfet de la Ville, à son épouse incomparable et pour lui-même a fait (ce monument).}
\comm{Le texte est gravé en colonne étroite, une ou deux syllabes par ligne, disposition que les deux bases reproduisent. Il ne contient aucune date : l’année 351 de la base dépend de l’identification de Clodius Adelfius à Clodius Celsinus Adelphius, préfet de la Ville en 351, et marque un terme postérieur, non l’année du décès (chapitre~\ref{ch:fondement}). La bibliographie fait de ce préfet l’époux de la poétesse chrétienne Faltonia Betitia Proba ; l’identification et cette relation sont à contrôler dans la \emph{Prosopography of the Later Roman Empire}. Le texte, sans formule chrétienne, est classé dans le corpus en raison de sa publication parmi les inscriptions chrétiennes (ICVR I 19).}

\dossier{38}{Privata, \emph{in pace} et le monogramme}
\trad{Privata, qui vécut cinquante-trois ans, huit mois et deux jours, a été déposée le quatrième jour avant les nones de février (2 février) ; bien méritante, en paix ☧ ; sous le consulat de Decentius et de Paulus (352).}
\comm{Dossier de référence pour la question du monogramme (chapitre~\ref{ch:dieu_christ}). L’EDCS développe \emph{Chr(isti)} ; l’EDB transcrit \emph{((monogramma Christi))} ; l’estampage conservé à l’Ohio State University sous le numéro ILCV 3252A montre un christogramme placé après IN PACE. La pierre porte donc un signe, non le mot écrit en lettres. La traduction rend le signe par ☧ sans décider s’il se lit \emph{Christi} (« dans la paix du Christ ») ou s’il est juxtaposé à la formule. Le texte provient, selon l’EDB, du cimetière à ciel ouvert proche de la basilique Saint-Pierre ; il est conservé au Lapidario Cristiano des Musées du Vatican. La forme \emph{meses} pour \emph{menses} est signalée par les deux bases.}

\dossier{43}{Anastasius, poème de Constantia}
\trad{(Vers) Au nom [du Christ …]. Affligée, [Constantia écrit ces vers] pour Anastasius, son époux, qui [a échangé] la lumière [contre d’amères ténèbres] ; sa vie [s’arrêta] à quatre fois dix [et cinq ans] ; combien vite il fut ravi [à l’amour de son épouse chérie] ! Pleuré [quand Janus prenait] son douzième [lever], il fut enseveli [l’année que conduisait] le consul Arbitio (355).}
\comm{Les deux bases donnent la même restitution, empruntée à l’édition : chaque vers est conservé sur sa moitié gauche, et la seconde moitié est reconstituée selon le mètre et le sens. Le nom du Christ, dans l’invocation initiale, est entièrement restitué ; l’EDB décrit un chrisme avec alpha et oméga, qui rend l’invocation vraisemblable sans la rendre lisible. L’année repose sur le nom d’Arbitio, lu (\emph{Arbitio co[nsul]}) ; la date du jour (\emph{duodecimum [Ianus]}) est une restitution. L’âge, « quatre fois dix [et cinq] », est en partie restitué.}

\dossier{70}{Secunda, \emph{virgo fidelis}}
\trad{À Secunda, d’une bonté admirable, qui vécut d’une foi pure vingt ans ; chaste, elle a cessé (de vivre) en paix, vierge fidèle ; bien méritante, elle repose ; aux ides de juillet (15 juillet) ; colombe sans fiel ; sous (le consulat de) M(amertinus) et N(evitta) (362).}
\comm{L’un des deux seuls textes datés par les seules initiales des consuls (chapitre~\ref{ch:fondement}) ; l’identification de l’année 362 repose sur le couple M. et N., que seul le collège de Mamertinus et Nevitta satisfait dans la période. Le vocabulaire est celui de la virginité consacrée probable (\emph{pura fide}, \emph{pudica}, \emph{virgo fidelis}) ; l’image de la colombe sans fiel (\emph{palumbo sine felle}, au masculin) appartient au langage chrétien de l’innocence. L’EDB lit \emph{quiescet} et conserve \emph{\{id\}}, que l’EDCS considère également comme superflu.}

\dossier{72}{Flavia Balbilla : \emph{Dis Manibus} et \emph{in pace}}
\trad{Consacré aux dieux Mânes. Flavia Balbilla vécut vingt-huit ans, six mois et douze jours ; elle a reposé en paix sous le consulat de Mamertinus et de Nevitta, Maximus étant préfet de la Ville, le cinquième jour avant les calendes de février (28 janvier 362).}
\comm{Seul texte de l’ensemble C qui porte la dédicace aux Mânes, associée au verbe chrétien du repos (\emph{requievit in pace}). Les deux formules coexistent sans difficulté apparente pour le commanditaire ; la dédicace \emph{D. M. s.} est alors une formule d’ouverture traditionnelle plus qu’une profession religieuse. Le texte est aussi le seul du corpus à dater par le préfet de la Ville en plus des consuls (chapitre~\ref{ch:rangs}). L’EDB transmet la graphie \emph{Nebidda}, que l’EDCS normalise.}

\dossier{73}{Eleuteria : l’esprit, les saints et les hauteurs}
\trad{(Vers) [Voici] la demeure éternelle où tu reposes désormais en sécurité ; […] ton esprit, se retirant de la chair, [vit avec / est associé aux] saints, pour tes mérites et tes si grandes œuvres ; [toi qui …] tu reposes toujours en sécurité ; [tu as confié ton corps] à la terre et le gage reçu ; […] à la splendeur avec une claire lumière ; […] toi qui t’es toujours consacrée à Dieu ; […] mais pour tes actions tu es appelée vers les hauteurs ; […] Pacatianus doit beaucoup […]. Pacatianus a fait cet emplacement pour Eleuteria, sa mère, qui vécut environ soixante-quinze ans, deux mois et dix jours ; déposée le quatrième jour avant les calendes de juin (29 mai) ; elle dort [en paix], sous le consulat de notre seigneur Claudius Iulianus Auguste pour la quatrième fois et de Flavius Sallustius (363). Hesperius, qu’elle a nourri, l’a gravé.}
\comm{Le poème est conservé sur sa partie droite, et les deux bases restituent différemment le début des vers (chapitre~\ref{ch:typologie}). La traduction ne retient que ce qui est commun : l’esprit se sépare de la chair (lu), il est en relation avec les saints (\emph{sanctis}, lu ; le verbe est restitué), la défunte s’est consacrée à Dieu (\emph{deo dica⟨vi⟩sti}, lu dans l’EDB, \emph{dicasti} dans l’EDCS) et elle est appelée « vers les hauteurs » (lu). Le séjour \emph{in caelo} et la crainte de Dieu (\emph{[deu]m metuisti}) ne sont que des restitutions de l’EDCS ; l’EDB lit \emph{[- - - de]metuisti}. Le texte nomme trois personnes : la défunte, son fils Pacatianus, qui fait l’emplacement, et Hesperius, son nourricier, qui grave. L’EDB décrit des palmes, un soleil, un chrisme, une croix gammée, un alpha et un oméga, disposés à la fin des vers.}

\dossier{110}{Benenatus, ancien \emph{comes}, et Theophila}
\trad{(Grec) En paix a reposé Benenatus, ancien \emph{comes}, ayant vécu soixante-quinze ans ; il gît en ce lieu, sorti de la vie le quatrième jour avant les ides de mars (12 mars), sous le consulat de Lupicinus et de Iovinus (367) ; moi, Theophila, j’ai fait (ce monument) pour mon époux. — (Latin) À Benenatus, Theophila, son épouse, a fait (ce monument).}
\comm{Épitaphe bilingue, dont le grec porte l’essentiel (chapitre~\ref{ch:grec}). Le formulaire grec (\emph{ἐν εἰρήνῃ ἀνεπάη}) est l’équivalent du \emph{requiescit in pace} latin ; la date romaine est transposée en grec. L’abréviation \emph{ἀπ(ὸ) κο(μήτων)} désigne un ancien \emph{comes}, sans préciser la fonction. Les graphies \emph{κῖτε} et \emph{ἐξελτών} reflètent la prononciation contemporaine. Le nom de l’épouse, \emph{Theophila}, est théophore.}

\dossier{118}{Ianuaria : l’âge à l’heure près}
\trad{Sous le consulat de Lupicinus et de Iovinus (367), Ianuaria, qui vécut vingt-quatre ans, un jour et deux heures et demie, s’en est allée un mardi.}
\comm{Seul texte du corpus qui donne, sans restitution, une durée de vie en heures (chapitre~\ref{ch:personnes}). L’EDB ne transcrit pas \emph{s(emis)} et lit \emph{XXI III}, ce qui pourrait signaler une disposition particulière des chiffres sur la pierre ; elle décrit l’image de la défunte et des colombes. Le quantième n’est pas donné : le mardi ne peut être contrôlé (chapitre~\ref{ch:calendrier}). Les deux noms consulaires sont gravés au datif ou à l’ablatif avec \emph{in consulatu}, construction que l’EDCS signale par \emph{(!)}.}

\dossier{126}{Un fils né en 350, mort en 368}
\trad{Ses parents affligés ont fait ce que lui-même aurait dû faire ; malheureux parents, (contraints) par la douleur à ce devoir, ils ont fait aussi ce tombeau. Il naquit sous le consulat de Flavius Sergius et de Flavius Nigrinianus (350), le troisième jour avant les ides d’août (11 août), un samedi ; il a été déposé dans cette demeure éternelle sous le consulat de nos seigneurs Valentinien et Valens, Augustes pour la deuxième fois (368), le septième jour avant les ides de décembre (7 décembre), un samedi ; il vécut dix-neuf ans, trois mois et vingt-six jours ; en paix.}
\comm{Le topos de l’inversion des devoirs (les parents font pour l’enfant ce que l’enfant aurait dû faire pour eux) ouvre le texte. Les deux dates sont complètes. Le calcul confirme le samedi de la naissance ; le 7 décembre 368 est un dimanche, soit un écart d’un jour avec le samedi gravé. Les mois et les jours de l’âge concordent exactement avec l’intervalle entre les deux dates, mais l’intervalle est de dix-huit ans et non de dix-neuf (chapitre~\ref{ch:calendrier}). La durée gravée, \emph{XVIIII}, est lue identiquement par les deux bases ; l’écart d’une année reste inexpliqué. La tombe est une \emph{domus aeterna}.}

\dossier{133}{Felix, « appelé à la sainte foi »}
\trad{Felix, appelé à la sainte foi, s’en est allé en paix ; un si grand amour et une si grande affection sont gardés pour lui par ses amis à jamais ; de son vivant, il fut un réconfort, miséricordieux, connu de tous. Agrippina a fait (ce monument) pour son très doux mari, avec qui elle vécut sans blessure de l’âme trois ans et dix mois ; il fut dans le siècle, tant qu’il vécut, trente-deux ans ; déposé le treizième jour avant les calendes de septembre (20 août), sous le consulat de Valentinien, très noble enfant, et de Victor (369).}
\comm{L’expression \emph{sanctae fidei vocitus} peut évoquer une conversion ou un baptême ; le texte ne précise pas lequel (chapitre~\ref{ch:foi_eglise}). La durée de la vie commune (trois ans et dix mois) avait été enregistrée comme un âge dans la table des personnes ; elle est corrigée (chapitre~\ref{ch:personnes}). L’éloge de Felix (\emph{solacius}, \emph{misericors}) relève de la charité ; l’expression \emph{fuit in saeculum} oppose la vie terrestre à ce qui suit. L’EDB rattache le texte à la basilique Saint-Pierre.}

\dossier{134}{Une dalle familiale : un prêtre, un diacre et leur sœur}
\trad{(a) Endormissement de Dulcitius, prêtre, aux calendes de février (1\textsuperscript{er} février), et de Gregorius, diacre, le treizième jour avant les calendes de mai (19 avril), et de Catafronia (EDB : Califronia), leur sœur, le cinquième jour avant les ides de mars (11 mars). (b) Déposition de Faustina le quatrième jour avant les calendes de novembre (29 octobre), sous le consulat de Valentinien, très noble enfant, et de Victor (369). (c) Endormissement d’Agapè le quatrième jour avant les nones (EDB : aux nones) d’octobre, en paix.}
\comm{La grande dalle (105 × 190 × 15 cm selon l’EDB), retrouvée en morceaux remployés dans un mur de la basilique des Saints-Nérée-et-Achillée, porte trois textes distincts, que l’EDB distingue (a, b, c). Seule la partie b est datée : l’année de classement ne vaut pas pour les autres défunts (chapitre~\ref{ch:evenements}). Le texte atteste une famille cléricale (un prêtre, un diacre et leur sœur) et l’emploi de \emph{dormitio} pour désigner la mort. Les deux bases divergent sur le nom de la sœur et sur la date d’Agapè.}

\dossier{140}{Bonifatia, \emph{in pace Chr(isti)}}
\trad{À Bonifatia, bien méritante, dans la paix du Christ, qui vécut vingt-huit ans ; déposée en paix le douzième jour avant les calendes de septembre (21 août), sous le consulat de Valentinien et de Valens pour la troisième fois (370).}
\comm{La formule \emph{in pace Chr(isti)} suit \emph{bene merenti} et précède une seconde mention de la paix (\emph{deposita in pace}). Les deux bases transcrivent une abréviation, mais l’EDB décrit un chrisme parmi les données matérielles : il est probable que l’abréviation développée rende un monogramme employé comme compendium (chapitre~\ref{ch:dieu_christ}). L’EDB rattache le texte au cimetière de Pontien.}

\dossier{151}{Muscula, dite aussi Galatea}
\trad{Ici gît Muscula, dite aussi Galatea, qui vécut deux ans, deux mois et dix-sept jours ; déposée le quinzième jour avant les calendes d’août (18 juillet), sous le consulat de Gratien Auguste pour la deuxième fois et de Probus (371), dans la paix du Christ.}
\comm{La formule \emph{in pace Chr(isti)} occupe une place inhabituelle : elle clôt le texte, après la date consulaire. Comme au \nnum{140}, les deux bases transcrivent une abréviation et l’EDB décrit un chrisme sur le support (chapitre~\ref{ch:dieu_christ}). Le double nom (\emph{quae et}) n’a pas de portée religieuse. Le texte provient, selon l’EDB, du cimetière de Commodilla ; il est conservé au Lapidario Cristiano.}

\dossier{183}{Petronia, vierge, de la neuvième région}
\trad{Déposition de Petronia le sixième jour avant les calendes de décembre (26 novembre), après le consulat de Gratien Auguste pour la troisième fois et de Flavius Equitius (375) ; elle vécut vierge environ dix-neuf ans ; elle est en paix ; (elle était) de la neuvième région, du côté du lac du \emph{Cuniculus}.}
\comm{Le mot \emph{virgo} peut désigner une jeune fille non mariée ou une vierge consacrée ; rien dans le texte ne permet de choisir (chapitre~\ref{ch:foi_eglise}). L’indication de domicile dans la neuvième région augustéenne, \emph{a lacu cuniculi}, est exceptionnelle dans le corpus ; la localisation de ce repère relève de la topographie de Rome et n’est pas proposée ici. L’EDB rattache le texte au cimetière de Zénon et décrit un chrisme, un paon et un rameau d’olivier ; la plaque est conservée au Collegio Teutonico.}

\dossier{202}{Cinnamius Opas, lecteur du titre de Fasciola}
\trad{Cinnamius Opas, lecteur du titre de Fasciola, ami des pauvres, qui vécut quarante-six ans, sept mois et neuf jours ; déposé en paix le dixième jour avant les calendes de mars (20 février), sous le consulat de Gratien pour la quatrième fois et de Merobaudes (377).}
\comm{L’un des trois lecteurs rattachés à un titre de la Ville (chapitre~\ref{ch:rangs}). Le \emph{titulus Fasciolae} est de nouveau nommé au \nnum{421} (398), avec le lecteur Pascentius. L’âge du défunt montre que la fonction de lecteur pouvait être exercée à l’âge mûr. La qualification \emph{amicus pauperum} associe le clerc à l’assistance. Le texte n’est pas apparié à l’EDB.}

\dossier{209}{Victorinus achète son emplacement au fossoyeur Martinianus}
\trad{Cet emplacement est celui de Victorinus, qui l’a acheté au fossoyeur Martinianus, sous (le consulat de) Valens pour la sixième fois et de Valentinien (pour la deuxième fois) (378).}
\comm{Le plus ancien acte d’achat daté du corpus qui nomme le vendeur, un \emph{fossor} (chapitre~\ref{ch:sepultures}). L’acte, sans commémoration, n’entre pas dans l’ensemble C. L’EDB restitue l’itération du second consul (\emph{Valentin[iano bis]}) et décrit, en tête du texte, un monogramme du Christ avec alpha et oméga : le signe chrétien marque ici un document de propriété.}

\dossier{214}{\emph{Ancilla dei} : « prie pour ce seul enfant »}
\trad{Ici repose une servante de Dieu, qui de ses biens a possédé tout entière cette demeure, que ses amies pleurent et dont elles cherchent le réconfort. Prie pour ce seul (enfant), la descendance que tu as laissée survivante. — (EDB seule :) Tu demeureras dans le repos éternel de la félicité ; le [quantième anomal] jour avant les calendes d’octobre. Cucurbitinus et Abundantius reposent ici ensemble. Sous le consulat de nos seigneurs Gratien pour la cinquième fois et Théodose, Augustes (380).}
\comm{Le seul témoignage assuré du corpus d’une demande d’intercession adressée au défunt (chapitre~\ref{ch:mort_esperance}). La défunte, \emph{ancilla dei}, est priée d’intercéder pour l’enfant qui lui survit. L’expression \emph{de sua omnia possedit domum ista} indique qu’elle a acquis la tombe sur ses propres biens. L’EDCS s’arrête à \emph{reliquisti} ; l’EDB transmet trois lignes supplémentaires, qui donnent la date et nomment deux autres défunts inhumés dans la même tombe. L’année 380 repose sur ces lignes. La plaque (69 × 128 cm) est conservée à la villa Albani ; l’EDB décrit des oiseaux tenant une couronne.}

\dossier{231}{Celerinus, enfant \emph{fidelis}}
\trad{Innocent, il s’en est allé, (enfant) des innocents, miséricordieux il fut. Qui n’a pas souffert de ton âge et versé de pieuses larmes ? En toi était attendue l’espérance de l’avenir, par toi une gloire durable, Celerinus, mon fils. Fidèle, tu reposes en paix. Il vécut un an et quatre mois ; déposé le sixième jour avant les calendes d’octobre (26 septembre), sous (le consulat d’)Antonius et de Syagrius (382).}
\comm{L’application de \emph{fidelis} à un enfant de seize mois suggère qu’il était baptisé (chapitre~\ref{ch:foi_eglise}). Le texte développe une rhétorique de la perte : l’enfant portait l’espérance et la gloire de la famille. L’EDB lit \emph{d(iebus) VI} là où l’EDCS lit \emph{d(epositus) VI} : la seconde lecture fait de \emph{VI Kal. Oct.} la date de la déposition, la première l’intègre à l’âge. La traduction suit l’EDCS.}

\dossier{232}{\emph{Intra limina sanctorum} : une sépulture recherchée}
\trad{[…]na, servante dans le Seigneur (?), […] a nourri […] les siens ; [amie des] pauvres ; [qui pour de si grands] mérites a reçu [une sépulture à l’intérieur des] seuils des saints ; […] a reçu [ce que beaucoup désir]ent et que peu obtiennent ; [sous le consulat d’Antonius et de] Syagrius (382).}
\comm{Le texte est conservé sur sa partie droite. Trois éléments intéressent l’histoire religieuse (chapitres~\ref{ch:mort_esperance} et~\ref{ch:sepultures}). (1) \emph{in dom(ino?) cultrix} : l’abréviation est développée avec un point d’interrogation par l’éditeur ; le texte n’est retenu que dans le profil large. (2) \emph{[- - - l]imina sanctorum} : la fin de l’expression est lue ; le mot qui la précédait est restitué. (3) \emph{[- - -]n(t) et rari accipiun(t)} : la seconde moitié de l’antithèse est lue. Le texte présente la sépulture près des saints comme la récompense des mérites d’une femme charitable et comme un privilège rare. Il ne dit pas quels saints, ni où ; il est d’origine incertaine.}

\dossier{244}{Proiecta : une épigramme de Damase}
\trad{(Vers) Que dire ? Ou me taire ? La douleur même m’interdit de parler. Ce tombeau retient les larmes — sache-le — des parents de Proiecta, qui avait été unie à Primus, son mari, belle de sa grâce, contente de sa seule pudeur. Hélas, trop chérie de l’amour d’une mère malheureuse ! Apprends ce qu’elle eut de plus que bien d’autres : après les premiers liens du mariage, arrachée aux yeux de Florus, son père, elle s’en est allée, désirant monter vers la lumière éthérée du ciel. Ces consolations des pleurs, Damase les offre à tous. — Elle vécut seize ans, [dix mois] et vingt-cinq jours ; déposée le troisième jour avant les calendes de janvier (30 décembre), sous le consulat de Flavius Merobaudes et de Flavius Saturninus (383).}
\comm{L’épigramme est signée par l’évêque de Rome Damase (366–384), et la formule consulaire la date de la dernière année de son épiscopat. Elle figure dans l’édition commentée de Trout (\emph{Damasus of Rome. The Epigraphic Poetry}, 2015, n\textsuperscript{o}~51 selon la concordance de l’EDCS) et dans les \emph{Carmina Latina Epigraphica} (CLE 670). Le texte ne nomme ni le Christ ni Dieu ; l’espérance de la défunte est exprimée dans le langage de l’ascension céleste. Le rapprochement avec la Proiecta nommée sur le coffret du trésor de l’Esquilin a été discuté dans la bibliographie consacrée à ce trésor ; il n’est pas repris ici. Les mots \emph{m(enses) X} sont signalés par l’EDCS comme une addition (\emph{«  »}).}

\dossier{249}{« Qui a reçu l’Esprit saint »}
\trad{[…] qui a reçu l’[Esprit] saint, [le seizième jour avant les calendes] de mai (16 avril ?), sous le consulat de Merobaudes [pour la deuxième fois] (383).}
\comm{Dossier examiné au chapitre~\ref{ch:esprit_trinite}. Sont lus \emph{-tus est qui ac-}, \emph{-tum spiritum}, \emph{Mai(as) Merobaud-}. Le verbe (\emph{ac[cepti]} dans l’EDCS, \emph{ac[cepit]} dans l’EDB), le début de l’adjectif, le quantième et l’itération du consulat sont restitués. La mention de l’Esprit saint est probable ; le rite qu’elle désigne et sa date ne peuvent être précisés. Le fragment (20 × 23,5 cm) provient, selon l’EDB, du cimetière de Nicomède, sur la via Nomentana.}

\dossier{268}{Lucifer, par son fils Benedictus}
\trad{Ce tombeau de son père, son fils Benedictus a voulu qu’il fût fait, par amour et en souvenir de son père, pour Lucifer, homme très innocent, [ami] des amis, qui, tant qu’il fut dans les choses humaines, parut très disponible envers tous et doux ; car il vécut toujours sans [faute de son] temps soixante-cinq ans, six mois et douze jours ; il mourut [en paix (?)]. — (EDB seule :) Le septième jour avant les calendes de juillet (25 juin), sous le consulat de Richomeres et de Clearchus (384).}
\comm{L’EDCS s’arrête à \emph{decess[it]} ; l’EDB transmet la dernière ligne, qui donne la date et fonde l’année 384 (chapitre~\ref{ch:fondement}). L’EDB souligne de nombreuses lettres, lues par les éditeurs anciens et aujourd’hui perdues. Le nom du défunt, \emph{Lucifer}, et celui du fils, \emph{Benedictus}, appartiennent l’un à l’onomastique traditionnelle, l’autre à l’onomastique chrétienne (chapitre~\ref{ch:noms}). La plaque (64 × 98 cm) est conservée à la Galleria Lapidaria des Musées du Vatican.}

\dossier{275}{Quirillus, Arménien de Cappadoce}
\trad{Citoyen arménien de Cappadoce, du nom de Quirillus, pieux, aimé de tous ; il vécut environ soixante-dix ans ; il nous a précédés dans la paix le sixième jour avant les ides de mars (10 mars), après le consulat de Richomeres et de Clearchus (385), un lundi ; (tombe) double avec Hyperechia, sa compagne.}
\comm{L’origine est mise en tête du texte, à l’accusatif (\emph{civem Armeniacum Cappadocem}), construction que l’EDCS signale. Le 10 mars 385 est bien un lundi (chapitre~\ref{ch:calendrier}). L’expression \emph{praecessit ad pacem} présente la mort comme un départ qui précède celui des survivants. La tombe est un \emph{bisomus} préparé pour le couple. L’EDB décrit un monogramme avec alpha et oméga ; la plaque est conservée au Lapidario Cristiano.}

\dossier{287}{Le jour, la lune et le signe}
\trad{[…] Honorius, [très noble enfant, et d’Euodius], \emph{vir clarissimus}, consuls (386), [le dixième jour avant les calendes de septembre (23 août)], un dimanche, [à la … heure de la nuit], au douzième jour de la lune, sous le signe [du Capricorne] ; […] trente jours ; […] le septième ; […] repose en paix ; […].}
\comm{Dossier examiné au chapitre~\ref{ch:calendrier}. Le quantième est entièrement restitué ; le jour de la semaine, l’âge de la lune et le signe zodiacal, en partie conservés, sont compatibles avec la restitution, mais ne la confirment pas indépendamment, puisque l’éditeur a pu la déduire de ces mêmes données. Les deux bases donnent le même texte ; l’EDB rattache la plaque au cimetière de Commodilla.}

\dossier{306}{Maximilla et Lucceia}
\trad{(Vers) Maximilla, capable [de garder l’alliance] du Christ et son propos : la mort lui a enlevé tous les biens de la vie ; mais il vit mieux, celui qui espère la vie éternelle du Christ, s’il médite en son âme la foi unique donnée à tous ; alors la mort elle-même pourra mourir, et la vie sera perpétuelle pour tout le temps ; il accordera à tous un pardon unique, mais selon les mérites. Toi, Lucceia, qui la soutiens, tu as fait (ce monument) au-dessus d’elle ; car, te souvenant d’elle quand tu acquittes tes vœux au tombeau, tu montres un esprit rempli d’une vie religieuse. — (Prose) D’une foi et d’une innocence admirables et d’un excellent propos, Maximilla, vierge, servante de Dieu, citoyenne de Pannonie, qui vécut cinquante ans, a été mise dans cette sépulture par Lucceia, \emph{clarissima femina}, fille de Viventius, \emph{vir} de mémoire clarissime, ancien préfet du prétoire et de la Ville éternelle, en raison de l’amitié qui les unissait ; elle a fini ses jours en même temps que sa mère Nunita, qui fut épouse d’un diacre et vécut soixante-dix ans ; le quatrième jour avant les nones d’octobre (4 octobre), sous le consulat des Flavii Timasius et Promotus, \emph{viri clarissimi} (389).}
\comm{Le texte le plus développé du corpus sur la vie éternelle (chapitres~\ref{ch:dieu_christ} et~\ref{ch:mort_esperance}). Il est gravé sur un sarcophage de marbre (72 × 212 cm) de l’aire à ciel ouvert de Saint-Sébastien. Le poème (lu dans les deux bases, sauf le début du premier vers) exprime une espérance construite : la vie éternelle du Christ, la mort de la mort, un pardon accordé selon les mérites. La partie en prose présente un milieu sénatorial (Lucceia, fille d’un ancien préfet) et clérical (Nunita, \emph{matrona diaconi}, selon la correction de l’EDCS). La ligne qui qualifie Maximilla de \emph{virgo ancilla dei civis Pannonia} manque dans le texte de l’EDB ; elle doit être contrôlée sur l’ICVR V 13355 avant toute exploitation. Le verbe \emph{[ve]getans} et la construction \emph{su[per h]anc} sont en partie restitués ; la traduction de la phrase qui les contient reste incertaine.}

\dossier{311}{Nisus : \emph{deposita in Chr(isto)}}
\trad{Déposé, Nisus, la veille des nones ; déposée dans le Christ le sixième jour avant les ides de février (8 février), sous le consulat de Maximus (388).}
\comm{Texte tracé dans le mortier frais, qui commémore deux défunts, dont le second n’est pas nommé dans la partie transmise. \emph{In Chr(isto)} est transcrit par les deux bases, mais l’EDB décrit un chrisme sur le support : l’abréviation développée rend probablement un monogramme (chapitre~\ref{ch:dieu_christ}). \emph{Maximo cons(ule)} désigne le consulat de Magnus Maximus, en usage à Rome au début de 388 ; l’EDB date de 389 et note que sa transcription diffère de l’ICVR. Datation contestée (chapitre~\ref{ch:fondement}).}

\dossier{321}{La colonne de Saint-Paul}
\trad{(c) Colonne de l’apôtre Paul ; jour de dédicace le [treizième] jour avant les calendes [de décembre (19 novembre)], sous le consulat de notre seigneur] Valentinien Auguste pour la quatrième fois et de Neoterius, \emph{vir clarissimus} (390), sous l’administration de Flavius Philippus, \emph{vir} [\emph{clarissimus}, …], le curateur étant Flavius Anastasius, [\emph{vir clarissimus}], tribun prétorien. — (a) Siricius, évêque, dévoué de tout son cœur. — (b) L’évêque dévoué.}
\comm{Seul texte non funéraire de l’ensemble B, inclus pour la datation consulaire et pour les noms de fonctions. L’EDB distingue trois textes gravés en trois endroits du fût : sous le chapiteau, au milieu de la colonne et la dédicace datée ; elle décrit un monogramme avec alpha et oméga après \emph{episcopus}. Le texte atteste l’engagement de l’évêque Sirice dans la construction de la nouvelle basilique de la via Ostiensis. Les noms de Flavius Philippus (fonctionnaire) et de Paul (l’apôtre) avaient été pris pour des noms consulaires par l’identification automatique (chapitre~\ref{ch:fondement}). Le quantième et une partie de la formule sont restitués.}

\dossier{324}{Hilaritas : fiançailles et mariage}
\trad{À Hilaritas, épouse bien méritante, qui vécut vingt-cinq ans ; fiancée à onze ans, elle fut fiancée sept ans ; elle vécut avec son mari sept ans et huit mois ; déposée le dix-septième jour avant les calendes de septembre (16 août), sous le consulat de Flavius Valentinien Auguste pour la quatrième fois et de Neoterius, \emph{vir clarissimus} (390).}
\comm{Le texte distingue l’âge aux fiançailles (onze ans), la durée des fiançailles (sept ans) et celle du mariage (sept ans et huit mois). Mariée vers dix-huit ans, Hilaritas est morte à vingt-cinq ans et huit mois environ ; l’âge gravé, vingt-cinq ans, est cohérent avec ce calcul (chapitre~\ref{ch:personnes}). Les deux bases concordent ; l’EDB rattache la plaque au cimetière de Commodilla.}

\dossier{352}{Eutychius, « croyant au Christ »}
\trad{(Vers) Ici a été placé un cruel trépas : dans sa première jeunesse, le jeune homme, arraché (à la vie), a franchi le seuil de la mort ; Eutychius, sage, sage, pieux et bienveillant, [qui], croyant [au Ch]rist, possède les récompenses de la lumière. Eutychius gît ici, lui-même l’unique honneur des siens. — Déposé le quatrième jour avant les nones d’octobre (4 octobre), sous le consulat de Théodose Auguste pour la troisième fois et d’Eugène (393) ; […] vingt-neuf ans et trois mois.}
\comm{Épitaphe métrique (CLE 1346), dont la pierre est perdue. Le nom du Christ est lu à partir de \emph{-ristum} ; la restitution \emph{[qui Ch]ristum}, contrainte par la finale, le sens et le mètre, est très probable, mais elle reste une restitution (chapitre~\ref{ch:typologie}). Le participe \emph{credens} construit avec l’accusatif exprime l’adhésion personnelle ; la récompense est « la lumière ». La date (\emph{Theodosio Aug. III et Eugenio}) est celle d’un consulat reconnu à Rome sous l’usurpation d’Eugène. Les graphies \emph{iubenta}, \emph{iubenis} et \emph{Eutichius}, normalisées par l’EDCS, sont conservées par l’EDB ; les deux bases corrigent \emph{benignus}, gravé avec un C.}

\dossier{378}{Une généalogie cléricale transmise par un seul témoin}
\trad{[Limenius] vécut […] ans ; […] Adeodata (?) […] ; […] il est tenu pour sage (?) […] ; […] doux pour moi d’abord […] ; […] il repose […] ; […] mois, à son épouse […] ; déposé le onzième jour avant les calendes d’octobre (21 septembre), [sous le consulat d’Olybrius et de Pro]binus (395). [Sa fille Vict]oria, déposée aux ides de janvier (13 janvier), [avec son mari …]orius. Telle est la lignée de la fille : Limenius, son père, prêtre ; petit-fils de Crescentius, petit-fils de […], diacre de la cinquième région ; petit-fils de […], arrière-petit-fils de l’évêque Marcellus, petit-fils de Crescentius, \emph{matrona} d’un prêtre ; arrière-arrière-petit-fils de Basilius ; belle-mère de l’évêque […], évêque Ursinus d’Albano.}
\comm{Ce texte, transmis par l’EDCS sous les références ICVR III 8161 et 8345, pose un problème de tradition (chapitre~\ref{ch:evenements}). Les fiches de l’EDB qui se rapportent au monument (ICVR III 8345, 8274c et 8387v) ne contiennent ni la formule consulaire ni la généalogie ; l’EDB date l’ensemble « 390–424 ». L’année 395 et toute la généalogie reposent donc sur un seul témoin, dont plusieurs termes sont restitués. La traduction suit l’ordre du texte transmis sans prétendre en reconstituer la syntaxe : l’accord des termes de parenté (\emph{nepos}, \emph{pronepos}, \emph{abnepos}, \emph{socrus}) et leur rattachement à la fille ou à d’autres personnes restent incertains. Si la lecture est confirmée, le texte atteste une famille où se succèdent prêtres, diacres et évêques, et le mariage des clercs. Contrôle de l’ICVR III 8345 indispensable.}

\dossier{382}{Panthagathè (?), mère de trois enfants}
\trad{(Vers acrostiches) [Avec la pal]me, elle possède les royaumes célestes […] ; [mais] d’une chasteté entière […] ; naguère, une si grande fidélité en toi […] ; une si grande pudeur, l’esprit […] ; cette âme reconnaissante […] ; […] l’enfant né avant terme, combien […] ; épouse généreuse pour moi et […] ; dorées sont les (demeures ?) célestes que toi, au ciel […] ; mère de trois enfants, [elle repose dans le sein de Dieu] ; son nom, [les initiales] des vers [le déclarent]. Elle vécut vingt-cinq ans ; [déposée] le dix-septième jour avant les calendes de septembre (16 août), sous le consulat d’Olybrius et de Probinus (395).}
\comm{Le poème est conservé sur sa moitié gauche ; les initiales des vers formaient le nom de la défunte, qui ne peut être établi sur les lettres conservées. L’EDCS ajoute, sur une autre face ou une autre partie du support, \emph{[Pa]nthagate}, absent de l’EDB. Le ciel (\emph{celestia regna}, \emph{caelestis}, \emph{caelo}) et la fidélité (\emph{fides}) sont lus ; Dieu n’apparaît que dans la restitution \emph{in sin[u dei pausat]}, commune aux deux bases (chapitre~\ref{ch:typologie}). Le vers \emph{acerbus natus} suggère une mort liée à une naissance prématurée, sans que la lacune permette de l’affirmer.}

\dossier{387}{Un achat sous l’épiscopat de Sirice ?}
\trad{[…] déposée le dix […], [sous le consulat de notre seigneur Honorius Auguste pour la troisième fois et] de notre seigneur Flavius Arcadius Auguste pour la quatrième fois (396) ; […] Eros, Leo et Setema [ont acheté] un emplacement [près du seuil de Pierre et de] Paul, sous l’épiscopat de Si[rice, aux prêtres Ursus] et Proclinus.}
\comm{Dossier examiné au chapitre~\ref{ch:sepultures}. Sont lus \emph{locu(m)}, \emph{Pauli sub episcopatu Si-} et \emph{et Proclino}. Le verbe d’achat, la mention de Pierre, la qualité de prêtres des vendeurs et le nom d’Ursus sont des restitutions. La datation épiscopale est vraisemblable (Sirice est évêque de 384 à 399) ; la localisation près de la tombe de Paul et l’intervention du clergé dans la vente ne le sont qu’au prix de ces restitutions. Le texte n’est pas apparié à l’EDB.}

\dossier{388}{Quodvultdeus}
\trad{Ici repose Quodvultdeus, homme d’honorable mémoire, qui vécut cinquante-sept ans ; déposé en paix le cinquième jour avant les ides d’octobre (11 octobre), sous le consulat — (EDB seule :) de nos seigneurs Arcadius Auguste pour la quatrième fois et Honorius Auguste pour la troisième fois (396).}
\comm{\emph{Quodvultdeus}, « ce que Dieu veut », est un anthroponyme : il contient le nom de Dieu sans constituer une invocation (chapitre~\ref{ch:noms}). L’EDCS s’arrête à \emph{conss(ulibus)} ; l’EDB transmet les noms consulaires, qui fondent l’année 396 (chapitre~\ref{ch:fondement}). La plaque (71 × 145 cm) est conservée au Lapidario Cristiano ; l’EDB décrit une feuille.}

\dossier{393}{Deux fragments d’une date consulaire}
\trad{[Sous le consulat d’Arcadius pour la] quatrième fois [et d’]H[onorius pour la troisième fois, Augustes] (396).}
\comm{L’EDCS publie sous deux numéros (anciens \nnums{393} et 394) des fragments que l’EDB réunit sous une fiche unique (ICVR IV 11126 a-b ; 16 × 13 × 2 cm). L’EDB ne lit que \emph{…IIII H…}, \emph{…ADIO…} et \emph{…o III Aug(usto)…}, compatibles avec le consulat d’Arcadius IV et Honorius III. Les deux fragments sont réunis sous ce numéro, et un seul texte daté de 396 est compté (chapitre~\ref{ch:corpus}). L’année repose sur une restitution guidée par les itérations (niveau R).}

\dossier{400}{Severianus : « dans la lumière du Seigneur »}
\trad{Doux et innocent, ici dort Severianus dans le sommeil de la paix ; il vécut environ cinquante ans ; son esprit a été accueilli dans la lumière du Seigneur le [quantième anomal] jour avant les calendes de mars, sous le consulat de nos seigneurs Flavius Caesarius et Nonius Atticus, \emph{viri clarissimi} (397) ; cet emplacement, Viricunda, son épouse, l’a acheté — Pascasius, \emph{ep(is)c(opus)} (?).}
\comm{Texte mixte (épitaphe et acte d’achat ; chapitre~\ref{ch:sepultures}). L’énoncé religieux est lu dans les deux bases : le défunt « dort dans le sommeil de la paix » et son esprit a été « accueilli dans la lumière du Seigneur ». L’EDB décrit, dans le texte même, un chrisme après le nom du défunt et une croix monogrammatique après \emph{susceptus est}, et, à la fin, un oiseau portant un rameau d’olivier et une croix monogrammatique dans une couronne. Le quantième \emph{VIIIIX} est anomal ; l’EDB reprend la note de Diehl, qui propose \emph{prid. Id. Febr.} et observe qu’aucun évêque de Rome ne porte le nom de Pascasius. L’EDB ne développe pas \emph{EPC}. La fonction de ce Pascasius dans l’acquisition reste indéterminée, et la traduction laisse la construction ouverte. L’âge, \emph{plus minus L}, ne comporte pas de jours : les « dix-sept jours » de la table des personnes provenaient de la date (chapitre~\ref{ch:personnes}).}

\dossier{405}{Un catéchumène}
\trad{À Bonifa(tius ?), qui repose bien, catéchumène, qui vécut environ soixante ans ; déposé le quatrième jour avant les calendes de juillet (28 juin), sous le consulat de Caesarius et d’Atticus (397).}
\comm{Seul texte du corpus qui nomme un catéchumène (chapitre~\ref{ch:foi_eglise}). Le défunt, mort à environ soixante ans, n’est pas dit baptisé ; le statut est gravé comme une qualification. Le nom est transmis sous la forme \emph{Bonifae}, que l’on peut comprendre comme une abréviation ou une graphie de \emph{Bonifatius}. L’EDB transcrit une partie effacée (\emph{〚- - -IV- - -〛}) et une addition (\emph{《d(epositus)》}) qui pourraient signaler une correction sur la pierre ; la plaque est publiée avec une reproduction dans un catalogue de 2000.}

\dossier{411}{Ursa : un prix d’emplacement}
\trad{Déposée en paix la jeune fille du nom d’Ursa, âgée de vingt-deux ans, sous le consulat de Flavius Caesarius et de Maximus (?), \emph{viri clarissimi} ; emplacement acheté (pour) trois (?) et onze sous.}
\comm{Texte mixte. L’âge est \emph{annorum viginti duo}, vingt-deux ans, et non vingt comme le portait la table des personnes ; \emph{puella} désigne ici une jeune femme (chapitre~\ref{ch:personnes}). Le collège transmis, \emph{Fl(avi) Caesari et Maximi}, ne correspond pas aux fastes de 397 (Caesarius et Atticus) : l’année est conservée sous réserve. Le prix (\emph{(t)res et s(olidis) XI} dans l’EDCS, \emph{⟨t⟩ris et s(olidis) (undecim)} dans l’EDB) est le seul du corpus ; sa construction reste incertaine (chapitre~\ref{ch:sepultures}).}

\dossier{433}{Un petit enfant parmi les astres}
\trad{(Vers) Ici gît, éteint au premier [seuil de la vie], un enfant très aimé [pendant un temps] bien court ; et je ne pense pas [qu’il soit permis d’accabler de] larmes [ce bienheureux], qui, dépouillé des liens du corps, [se réjouit parmi les astres], et n’a pas ressenti les [contagions] mauvaises [du sens] terrestre ; car après [avoir ajouté] neuf mois [à une année], il a accompli six jours de vie (de plus) [et s’en est allé vers les astres]. — Déposé sous le consulat de Lucius Mallius Theodorus, \emph{vir clarissimus} (399), la veille des nones […] ; il vécut un an, neuf mois [et six jours].}
\comm{Les vers sont conservés sur leur moitié gauche, et l’EDB souligne de nombreuses lettres de cette moitié, lues par les éditeurs anciens et aujourd’hui perdues. Les mots \emph{astris} et \emph{astra}, qui portent l’image de l’ascension, appartiennent entièrement aux fins de vers restituées (chapitre~\ref{ch:mort_esperance}). Ce qui est lu (\emph{corporis exutus vinclis}, \emph{nec mala terreni sensit}) suffit cependant à situer le poème dans le registre de la consolation, qui présente la mort de l’enfant comme une libération. L’EDB lit \emph{angusto} (« bref ») là où l’EDCS lit \emph{augusto}. L’année repose sur le consulat de Mallius Theodorus, lu ; le prénom \emph{Lucius} gravé est inhabituel pour ce consul.}

\dossier{447}{Une imprécation grecque}
\trad{(Grec ; partie transmise par l’EDB seule :) [Ici] repose Eu[…], [citoyen de Théo]dosiopolis, [ayant vécu] un peu plus de […] ans, [il s’est endormi] le quinzième jour avant les calendes de septembre (18 août), [sous le consulat de Vincentius] et de Fravitta, [clarissimes] (401). — (Partie commune aux deux bases :) Je vous adjure par [Dieu le Tout-]puissant : si quelqu’un [ose déranger] ses os [ou …], qu’il soit passible [du jugement / du châtiment] […] au [redoutable] jour [de la résurrection (EDCS) / du jugement (EDB)].}
\comm{Dossier examiné aux chapitres~\ref{ch:grec} et~\ref{ch:mort_esperance}. L’EDCS ne transmet que les cinq dernières lignes ; l’EDB en donne dix, qui contiennent le nom, l’origine et la date, et suit une autre édition. Dans la formule d’adjuration, sont lus \emph{ὁρκίζομε δὲ ὑμᾶς}, \emph{-ντοκράτορα}, \emph{εἴ τις το-}, \emph{τὰ ὀστέα} (EDB : \emph{ᾡτέα}) \emph{αὐτοῦ}, \emph{ἔνοχος τῇ} et \emph{ἡμαίραν τῆς}. Les deux bases restituent différemment la sanction : \emph{τῇ [κρίσει] … τῆς [ἀναστάσεως(?)]} dans l’EDCS, \emph{τῇ [κολάσει τὴν φοβερ]ὰν ἡμαίραν τῆς [κρίσεως]} dans l’EDB. Le texte atteste une menace eschatologique ; il n’atteste ni la résurrection ni le jugement par leur nom. L’expression \emph{ἔνοχος … τῇ κρίσει}, si elle est retenue, rappelle Mt 5, 21–22.}

\dossier{458}{Praetiosa, \emph{ancilla dei et Christi}}
\trad{La veille des ides de juin (12 juin) a reposé Praetiosa, jeune fille, vierge, de douze ans seulement, servante de Dieu et du Christ, sous le consulat de Flavius Vincentius et de Fravitta, \emph{viri clarissimi} (401).}
\comm{La formule \emph{ancilla dei et Chr(ist)i} associe Dieu et le Christ dans une qualification de dévotion ; elle ne contient ni « Père » ni « Fils » (chapitre~\ref{ch:dieu_christ}). L’EDB transcrit \emph{Chr(ist)i} : le I final est gravé, et il s’agit d’une abréviation en lettres, non d’un monogramme. L’association de \emph{virgo} et d’\emph{ancilla dei} pour une jeune fille de douze ans peut signaler une consécration précoce ou un éloge de l’innocence ; le texte ne permet pas de choisir. Le nom est \emph{Praetiosa} dans l’EDB, \emph{Praetioso} dans l’EDCS. La plaque provient, selon l’EDB, de la crypte de saint Corneille, dans le complexe de Calliste.}

\dossier{470}{Felicitas, « appelée d’un surnom trompeur »}
\trad{(Vers) Ce qu’offrent un doux enfant, ce qu’offrent de chers gages, ce tombeau le contient, qui retient de petits membres. Elle a causé une douleur sans fin, Felicitas, enfermée ici, malheureuse, appelée d’un surnom trompeur ; car, à elle que le premier jour avait mise au monde, (furent données) quatre années, et, comme elle commençait malheureusement la cinquième, un jour noir l’emporta et la plongea dans une mort prématurée. Sa mère et son père composent ces vers pour le tombeau, afin qu’en même temps, pour qui les lit, l’image revienne au cœur, et que les yeux, souvent séchés de pleurs, se mouillent de nouveau ; ainsi l’amour se soigne, et les Mânes n’ont que faire des vers. — Déposée en paix aux calendes de juillet (1\textsuperscript{er} juillet), sous le consulat de notre seigneur Théodose Auguste et de Flavius Rumoridus, \emph{vir clarissimus} (403).}
\comm{Épitaphe métrique (CLE 682) gravée sur un sarcophage, en deux parties séparées, selon l’EDB, par l’image d’une jeune orante. Le poème est entièrement classique dans sa langue et ses motifs : le jeu sur le nom (\emph{Felicitas} / \emph{infelix}), le \emph{funus acerbum}, la fonction consolatrice des vers, et même les Mânes, qui « n’ont que faire des vers ». La partie en prose se termine par \emph{deposita in pace}. Le texte illustre la coexistence, dans un même monument chrétien, du langage poétique traditionnel et du formulaire chrétien (chapitre~\ref{ch:mort_esperance}).}

\dossier{474}{Quintilianus, \emph{confirmans trinitatem}}
\trad{Quintilianus, homme de Dieu, […], affermissant (ou professant) la Trinité, […], aimant la chasteté, […], rejetant le monde, […], repose avec […], qui vécut environ dix […] ans ; sous le consulat de Théodose Auguste et de R[umoridus, \emph{vir clarissimus}] (403).}
\comm{Dossier examiné au chapitre~\ref{ch:esprit_trinite}. Le texte, identique dans les deux bases, est conservé sur sa moitié gauche. Les trois participes forment un éloge ascétique ; \emph{confirmans trinitatem} qualifie la foi du défunt, sans nommer les trois personnes. L’âge commence par \emph{X} et se poursuit dans la lacune : il est d’au moins dix ans. L’année repose sur \emph{Theodosio Aug(usto)}, lu, et sur l’initiale \emph{R}. La pierre, aujourd’hui perdue, provient selon l’EDB du cimetière à ciel ouvert proche de la basilique Saint-Pierre.}

\dossier{475}{\emph{Dedit deus, deus tulit}}
\trad{[…] Dieu a donné, Dieu a repris ; […] des calendes de novembre, sous le consulat de Théodose et de Rumoridus, \emph{viri clarissimi} (403).}
\comm{Citation libre de Jb 1, 21 (Vulgate : \emph{Dominus dedit, Dominus abstulit}), employée comme consolation (chapitre~\ref{ch:dieu_christ}). La forme \emph{deus … tulit} peut dépendre d’une version latine antérieure à la Vulgate ou d’une citation de mémoire. Le texte est un fragment ; l’EDB le rattache au cimetière d’Octavilla, près de Saint-Pancrace.}

